% ============================================================================
%  crystal_density_vacancy_diffraction.tex
%  VSEPR-SIM / thenewmethods
%
%  Document 1 of 3 -- The New Methods
%  Crystal structure, density, vacancy, and diffraction relations
%  Compact, invertible, system-grouped notation
% ============================================================================

\documentclass[11pt, letterpaper]{article}

\usepackage{amsmath, amssymb, amsthm}
\usepackage{geometry}
\usepackage{booktabs}
\usepackage{array}
\usepackage{xcolor}
\usepackage{microtype}
\usepackage{parskip}
\usepackage{hyperref}

\geometry{margin=1.1in}

\definecolor{accent}{RGB}{30, 80, 160}
\definecolor{note}{RGB}{120, 120, 120}

\hypersetup{colorlinks=true, linkcolor=accent, urlcolor=accent}

\title{\textbf{The New Methods I}\\[0.4em]
\large Crystal Structure \textbullet\ Density \textbullet\
Vacancy \textbullet\ Diffraction}
\author{VSEPR-SIM thenewmethods}
\date{}

% ── convenience macros ──────────────────────────────────────────────────────
\newcommand{\NA}{N_{\!A}}
\newcommand{\Qv}{Q_v}
\newcommand{\Nv}{N_v}
\newcommand{\dhkl}{d_{hkl}}
\newcommand{\argmin}{\operatorname{arg\,min}}

% ── section style ───────────────────────────────────────────────────────────
\usepackage{titlesec}
\titleformat{\section}{\large\bfseries\color{accent}}{}{0em}{}[\vspace{-0.3em}\rule{\linewidth}{0.6pt}]
\titleformat{\subsection}{\normalsize\bfseries}{}{0em}{}

% ============================================================================
\begin{document}
\maketitle
\thispagestyle{empty}

\noindent
Notation that is readable, compact, invertible where possible, and grouped
as systems where that helps.

% ────────────────────────────────────────────────────────────────────────────
\section{Crystal Structure Relations}
% ────────────────────────────────────────────────────────────────────────────

\subsection{FCC}

\[
\text{FCC:}\quad
\begin{cases}
  n = 4 \\[2pt]
  \mathrm{CN} = 12 \\[2pt]
  \mathrm{APF} = 0.74 \\[2pt]
  a = 2\sqrt{2}\,R \qquad \Longleftrightarrow \qquad R = \dfrac{\sqrt{2}}{4}\,a
\end{cases}
\]

\subsection{BCC}

\[
\text{BCC:}\quad
\begin{cases}
  n = 2 \\[2pt]
  \mathrm{CN} = 8 \\[2pt]
  \mathrm{APF} = 0.68 \\[2pt]
  a = \dfrac{4R}{\sqrt{3}} \qquad \Longleftrightarrow \qquad R = \dfrac{\sqrt{3}}{4}\,a
\end{cases}
\]

\subsection{HCP (for review, not default assumption)}

\[
\text{HCP:}\quad
\begin{cases}
  n = 6 \\[2pt]
  \mathrm{CN} = 12 \\[2pt]
  \mathrm{APF} = 0.74 \\[2pt]
  a = 2R
\end{cases}
\qquad \text{ideal ratio:}\quad \frac{c}{a} \approx 1.633
\]

% ────────────────────────────────────────────────────────────────────────────
\section{Density Relations}
% ────────────────────────────────────────────────────────────────────────────

Master equation for cubic unit cells:

\[
\boxed{\rho = \frac{n A}{\NA\, a^3}}
\]

\noindent
where $n$ = atoms per unit cell, $A$ = atomic weight, $\NA$ = Avogadro's number,
$a^3$ = unit cell volume.

\subsection{Inverse Forms}

Solve for atomic weight:
\[
A = \frac{\rho\, \NA\, a^3}{n}
\]

Solve for lattice parameter:
\[
a = \left(\frac{n A}{\rho\, \NA}\right)^{1/3}
\]

One line, invertible, no theatrical suffering.

% ────────────────────────────────────────────────────────────────────────────
\section{Vacancy Relations}
% ────────────────────────────────────────────────────────────────────────────

\subsection{Forward Form}

\[
\Nv = N\exp\!\left(-\frac{\Qv}{kT}\right)
\]

\subsection{Vacancy Fraction (often the more useful version)}

\[
\frac{\Nv}{N} = \exp\!\left(-\frac{\Qv}{kT}\right)
\]

\subsection{Inverse Form for $\Qv$}

Take the natural log:
\[
\ln\!\left(\frac{\Nv}{N}\right) = -\frac{\Qv}{kT}
\qquad\Longrightarrow\qquad
\Qv = -kT\ln\!\left(\frac{\Nv}{N}\right)
\]

\subsection{Inverse Form for Temperature}

\[
T = -\frac{\Qv}{k\ln\!\left(N_v/N\right)}
\]

\noindent
Memorization form: \emph{exponential forward, log backward}.

% ────────────────────────────────────────────────────────────────────────────
\section{Diffraction Relations}
% ────────────────────────────────────────────────────────────────────────────

\subsection{Bragg Law}

\[
n\lambda = 2d\sin\theta
\]

\subsection{Inverse Forms}

Solve for spacing:
\[
d = \frac{n\lambda}{2\sin\theta}
\]

Solve for angle:
\[
\theta = \sin^{-1}\!\left(\frac{n\lambda}{2d}\right)
\]

\noindent
When XRD reports $2\theta$ directly:
\[
\theta = \tfrac{1}{2}(2\theta)
\]

% ────────────────────────────────────────────────────────────────────────────
\section{Cubic Plane-Spacing Relations}
% ────────────────────────────────────────────────────────────────────────────

For cubic crystals:
\[
\dhkl = \frac{a}{\sqrt{h^2 + k^2 + l^2}}
\qquad\Longleftrightarrow\qquad
a = \dhkl\sqrt{h^2 + k^2 + l^2}
\]

\subsection{Combined Diffraction System}

Useful when solving structure problems:
\[
\begin{cases}
  n\lambda = 2\,\dhkl\,\sin\theta \\[4pt]
  \dhkl = \dfrac{a}{\sqrt{h^2+k^2+l^2}}
\end{cases}
\]

Substitute and solve for $a$:

\[
\boxed{a = \frac{n\lambda\,\sqrt{h^2+k^2+l^2}}{2\sin\theta}}
\]

That is the one-line usable relation to keep.

% ────────────────────────────────────────────────────────────────────────────
\section{Ionic Characterization}
% ────────────────────────────────────────────────────────────────────────────

Let $\Delta\chi = |\chi_A - \chi_B|$.

\[
\text{\%IC} = 100\!\left(1 - \exp\!\left[-\tfrac{1}{4}(\Delta\chi)^2\right]\right)
\]

System form (memorization form):
\[
\begin{cases}
  \Delta\chi = |\chi_A - \chi_B| \\[4pt]
  \text{\%IC} = 100\!\left(1 - \exp\!\left[-\dfrac{(\Delta\chi)^2}{4}\right]\right)
\end{cases}
\]

% ────────────────────────────────────────────────────────────────────────────
\section{Structure Classification Logic}
% ────────────────────────────────────────────────────────────────────────────

Input: $\{R,\,A,\,\rho,\,T,\,\theta,\,\lambda\}$.
Try FCC and BCC; for each compute $a(R)$, $\rho_{\mathrm{theoretical}}$,
$\dhkl$, and consistency with $T$.

Best structure via weighted residual:
\[
S^* = \argmin_{S\,\in\,\{\text{FCC},\text{BCC}\}}
\bigl[
  w_\rho\,|\rho_{\exp} - \rho_S|
  + w_d\,|d_{\exp} - d_S|
  + w_T\,\Pi_T(S)
\bigr]
\]

General classification objective:
\[
J(S) = \sum_i w_i\,|x_i^{\mathrm{measured}} - x_i(S)|
\qquad S^* = \argmin_S J(S)
\]

\noindent
Measured features: $x_i \in \{\rho,\,\dhkl,\,a,\,\mathrm{CN},\,\mathrm{APF},\,T\text{-stability}\}$.

% ────────────────────────────────────────────────────────────────────────────
\section{Ring Orbitals (compact Hückel form)}
% ────────────────────────────────────────────────────────────────────────────

For a cyclic $\pi$-system with $N$ p-orbitals:
\[
E_m = \alpha + 2\beta\cos\!\left(\frac{2\pi m}{N}\right),
\qquad m = 0, 1, \ldots, N-1
\]

Fill from lowest $E_m$ upward, 2 electrons per orbital.

\[
4n+2 \text{ pi electrons} \;\Rightarrow\; \text{aromatic}
\qquad
4n \text{ pi electrons} \;\Rightarrow\; \text{antiaromatic tendency}
\]

% ────────────────────────────────────────────────────────────────────────────
\section{Compact Memorization Block}
% ────────────────────────────────────────────────────────────────────────────

\begin{align*}
&\text{FCC: } a=2\sqrt{2}\,R,\quad n=4,\quad \mathrm{CN}=12,\quad \mathrm{APF}=0.74 \\
&\text{BCC: } a=\tfrac{4R}{\sqrt{3}},\quad n=2,\quad \mathrm{CN}=8,\quad \mathrm{APF}=0.68 \\
&\text{HCP: } a=2R,\quad n=6,\quad \mathrm{CN}=12,\quad \mathrm{APF}=0.74,\quad c/a\approx1.633 \\[4pt]
&\rho = \frac{nA}{\NA a^3} \\[4pt]
&\Nv = N\exp\!\left(-\frac{\Qv}{kT}\right)
\qquad
\Qv = -kT\ln\!\left(\frac{\Nv}{N}\right) \\[4pt]
&n\lambda = 2d\sin\theta
\qquad
d = \frac{a}{\sqrt{h^2+k^2+l^2}} \\[4pt]
&\Delta\chi = |\chi_A - \chi_B|
\qquad
\text{\%IC} = 100\!\left(1-e^{-(\Delta\chi)^2/4}\right) \\[4pt]
&S^* = \argmin_S J(S)
\end{align*}

% ────────────────────────────────────────────────────────────────────────────
\section{Full System Form}
% ────────────────────────────────────────────────────────────────────────────

\[
\begin{cases}
  \text{Given: } A,\,R,\,\rho_{\exp},\,\lambda,\,\theta,\,T,\,\chi_A,\,\chi_B \\[4pt]
  a_{\mathrm{FCC}} = 2\sqrt{2}\,R,\quad a_{\mathrm{BCC}} = \dfrac{4R}{\sqrt{3}} \\[6pt]
  \rho_{\mathrm{FCC}} = \dfrac{4A}{\NA\, a_{\mathrm{FCC}}^3},\quad
  \rho_{\mathrm{BCC}} = \dfrac{2A}{\NA\, a_{\mathrm{BCC}}^3} \\[6pt]
  d = \dfrac{n\lambda}{2\sin\theta} \\[6pt]
  \Delta\chi = |\chi_A - \chi_B| \\[4pt]
  \text{\%IC} = 100\!\left(1 - e^{-(\Delta\chi)^2/4}\right) \\[4pt]
  \dfrac{\Nv}{N} = \exp\!\left(-\dfrac{\Qv}{kT}\right) \\[6pt]
  S^* = \displaystyle\argmin_{S\in\{\text{FCC},\text{BCC}\}} J(S)
\end{cases}
\]

\end{document}
